\section{Pattern Planare Custom e Confronto con Zhang}
\label{sec:pattern_custom}

\subsection{Il principio geometrico comune}

Sia il metodo di Zhang sia CCalibGA sfruttano lo stesso principio:
\begin{quote}
\emph{Conoscere corrispondenze 2D--3D su un piano consente di stimare
la proiettività \(\mat{H}\) tra quel piano e il piano immagine;
imporre l'ortonormalità delle colonne di \(\mat{R}\) su \(\mat{H}\)
fornisce i vincoli lineari che identificano \(\mat{K}\).}
\end{quote}
La differenza tra i due approcci non risiede nella geometria della
calibrazione (identica), ma nel \emph{modo in cui le corrispondenze
vengono costruite}, come mostrato in Tabella~\ref{tab:confronto}.

\begin{table}[H]
\centering
\caption{Confronto geometrico tra Zhang e CCalibGA.}
\label{tab:confronto}
\begin{tabular}{p{3.2cm} p{5.0cm} p{5.0cm}}
\toprule
\textbf{Aspetto} &
\textbf{Zhang / scacchiera} &
\textbf{CCalibGA / custom} \\
\midrule
Struttura del target &
  Griglia regolare \(n{\times}m\) di quadrati bianco-neri &
  Qualsiasi immagine planare con texture ricca \\[4pt]
Estrazione dei punti &
  Angoli dei quadrati: posizione sub-pixel
  ricavata geometricamente dalla struttura &
  Keypoint affini rilevati sull'immagine di riferimento
  (\texttt{ccalib.CustomPattern}) \\[4pt]
Corrispondenze 3D &
  \emph{Implicite}: il punto \((i,j)\) ha coordinate
  \((id_x,\,jd_y,\,0)\) note a priori &
  \emph{Esplicite}: le coordinate 3D dei keypoint si
  leggono dalla geometria nota del riferimento \\[4pt]
Matching &
  Non necessario: la struttura metrica è già
  codificata nell'ordine della griglia &
  Descrittori affini + RANSAC~\cite{Fischler1981}:
  ogni keypoint del riferimento viene abbinato
  al corrispondente nella fotografia \\[4pt]
Robustezza &
  Alta su pattern ideali; sensibile a riflessi
  e occlusioni &
  RANSAC scarta outlier geometrici prima
  della stima di \(\mat{H}\) \\[4pt]
Geometria successiva &
  \multicolumn{2}{c}{Identica: DLT per \(\mat{H}\),
  vincoli IAC, Cholesky per \(\mat{K}\), LM per
  raffinamento} \\
\bottomrule
\end{tabular}
\end{table}

\subsection{Flusso di matching con pattern custom}

\begin{figure}[H]
\centering
\begin{tikzpicture}[>=Stealth,
  node distance=0.9cm and 3.0cm,
  rbox/.style={draw, rounded corners=4pt, fill=red!8,
               minimum width=3.4cm, minimum height=0.78cm,
               align=center, font=\small},
  bbox/.style={draw, rounded corners=4pt, fill=blue!8,
               minimum width=3.4cm, minimum height=0.78cm,
               align=center, font=\small},
  gbox/.style={draw, rounded corners=4pt, fill=green!10,
               minimum width=3.4cm, minimum height=0.78cm,
               align=center, font=\small},
  arr/.style={->, thick}]

  % Colonna sinistra: riferimento
  \node[rbox] (ref)   {Immagine di riferimento\\(pattern custom)};
  \node[rbox, below=of ref]  (kpR)  {Keypoint affini\\+ descrittori};
  \node[rbox, below=of kpR]  (p3D)  {Coordinate 3D\\
                                      \(\{(X_i,Y_i,0)\}\)};

  % Colonna destra: fotografia
  \node[bbox, right=of ref]  (foto) {Fotografia di\\calibrazione};
  \node[bbox, below=of foto] (kpF)  {Keypoint affini\\+ descrittori};
  \node[bbox, below=of kpF]  (p2D)  {Coordinate 2D\\
                                      \(\{(u_i,v_i)\}\)};

  % Blocchi centrali
  \node[gbox, below=1.8cm of kpR, xshift=2.6cm]
        (match) {Matching + RANSAC};
  \node[gbox, below=of match] (H) {DLT \(\Rightarrow\hat{\mat{H}}\)};
  \node[gbox, below=of H]     (K) {Vincoli IAC
                                    \(\Rightarrow\mat{K}\)};

  \draw[arr] (ref)  -- (kpR);
  \draw[arr] (kpR)  -- (p3D);
  \draw[arr] (p3D)  -- (match);

  \draw[arr] (foto) -- (kpF);
  \draw[arr] (kpF)  -- (p2D);
  \draw[arr] (p2D)  -- (match);

  \draw[arr] (match) -- (H);
  \draw[arr] (H)     -- (K);
\end{tikzpicture}
\caption{Flusso di costruzione delle corrispondenze con pattern custom.
         Il ramo rosso (riferimento) e quello blu (fotografia) si
         fondono nel blocco di matching.}
\label{fig:flusso_custom}
\end{figure}

\subsection{Proiezione del sistema di riferimento sul target}
\label{sub:assi}

Una volta stimata la coppia \((\mat{K},[\mat{R}_j|\mathbf{t}_j])\),
la validità geometrica si verifica visivamente proiettando il sistema
di riferimento 3D sull'immagine:
\[
  \homo{x}_\text{asse}
  = \mat{K}\,[\mat{R}_j\mid\mathbf{t}_j]\,\homo{X}_\text{asse}.
  \label{eq:assi_proj}
\]
I tre versori dell'origine del pattern (rosso \(=x\), verde \(=y\),
blu \(=z\)) si proiettano nel piano immagine: se la calibrazione è
corretta, le linee convergono coerentemente all'origine del target e
rispettano le proporzioni metriche dell'oggetto.